\subsection{SCBA iterations}
\label{sub:scba}
% ----------------------------------------------------------------------

\Cref{eq:system,eq:keldysh_eq,eq:sigma_phph,eq:sigma_r} together define
a nonlinear fixed-point problem $\Sigma_s = \mathcal{F}[G[\Sigma_s]]$,
solved iteratively in the SCBA. Convergence is two-fold: numerical
(consecutive iterates close in some norm) and physical (current
conservation, \cref{eq:I_local}). Mixing schemes are necessary because
pure fixed-point iteration of \cref{eq:sigma_phph} is often divergent or
oscillatory at finite anharmonicity~\cite{wang_nonequilibrium_2007,luisier_atomistic_2012}.

\subsubsection{Linear mixing}

The simplest scheme,
\begin{equation}\label{eq:linear_mix}
  \Sigma^{(n+1)} = (1-\alpha)\,\Sigma^{(n)} + \alpha\,\Sigma^{\mathrm{new}},
\end{equation}
with damping parameter $\alpha \in (0,1]$. Linear mixing converges if
the fixed-point Jacobian
$J = \partial\mathcal{F}/\partial\Sigma_s$ has spectral radius
% REVIEW: the linear-mixing map is M=(1-alpha)I+alpha J=I-alpha(I-J), so the
% convergence condition is sigma((1-alpha)I+alpha J)<1 (was sigma(I-alpha J)<1).
% The final bound alpha < 2/(1+|sigma(J)|_max) is correct.
$\sigma\bigl((1-\alpha) I + \alpha J\bigr) < 1$, i.e.\ $\alpha < 2/(1+|\sigma(J)|_{\max})$.
For weak anharmonicity ($|\sigma(J)|<1$) any $\alpha \le 1$ works; for
strong anharmonicity the spectrum of $J$ extends close to $-1$ in
some directions and small $\alpha\sim 0.1\text{--}0.3$ is required, with
linear convergence rate $\max_\lambda|1-\alpha+\alpha\lambda|$ set by the
eigenvalue of $J$ nearest $+1$ (which tends to $1$ as the iteration approaches
marginal stability).

\subsubsection{Anderson mixing}

Anderson acceleration~\cite{anderson_iterative_1965} constructs a quasi-Newton
update from a history of residuals,
\begin{equation}\label{eq:anderson}
  \vc{x}_{k+1}
  = \bigl(\vc{x}_k + \alpha\,\vc{f}_k\bigr)
    - \bigl(\Delta\vc{X} + \alpha\,\Delta\vc{F}\bigr)\,\vc{\gamma}_k,
\end{equation}
with $\vc{f}_k = G(\vc{x}_k) - \vc{x}_k$ the residual,
$\Delta\vc{X}, \Delta\vc{F}$ the matrices of consecutive iterate and
residual differences, and $\vc{\gamma}_k$ solving the regularized
least-squares problem
\begin{equation}
  (\Delta\vc{F}^H \Delta\vc{F} + \lambda \mathbb{I})\,\vc{\gamma}
  = \Delta\vc{F}^H \vc{f}_k,
  \qquad
  \lambda = 10^{-8}\,\tr(\Delta\vc{F}^H \Delta\vc{F}).
\end{equation}
For linear problems Anderson mixing is
essentially equivalent to GMRES applied to the residual
equation~\cite{walker_anderson_2011,rohwedder_analysis_2011}, with provable
% REVIEW: removed the explicit "q = 3c - c^2 for c < 2-sqrt(3)" -- it was
% internally inconsistent (2-sqrt(3) is the root of c^2-4c+1, unrelated to
% 3c-c^2; and 3c-c^2 > c for c in (0,2) is not an acceleration). The depth-1
% q-factor and its validity threshold are given explicitly by Bian-Chen-Kelley;
% reinsert the exact closed form from that reference if a number is wanted.
local $q$-linear convergence for depth-1
Anderson when the underlying iteration is a
contraction~\cite{toth_convergence_2015,bian_anderson_2021}. Local
convergence in nonlinear electronic-structure problems is well
established under mild Lipschitz assumptions on the
fixed-point map~\cite{rohwedder_analysis_2011,toth_convergence_2015}. In
practice, Anderson mixing dramatically accelerates convergence when the
fixed-point equation is well-conditioned, with rates approaching that
of Newton's method. However, storing the history is often prohibitive, so that
Anderson mixing becomes infeasible for large structures.

\subsubsection{Extrapolation and Newton root-finders for unstable fixed points}
\label{sub:rootfinders}

Linear and Anderson mixing both presuppose that the fixed point is
\emph{stable} under iteration, i.e.\ that the map Jacobian
$G'\equiv\partial\mathcal{F}/\partial\Sigma_s$ is contractive,
$|\sigma(G')|<1$. As the numerical broadening is taken to zero
(\cref{sub:eta_zero}) this fails for sharp, weakly damped resonances: a few
eigenvalues of $G'$ acquire modulus $\geq 1$ and no contraction-based scheme
converges. Three progressively stronger tools target this regime, all
MPI-distributed over the energy grid; their global reductions (the residual
inner products and norms, the RRE/RPM Gram solves $Z^{H}G'Z$, and the JFNK
trust-region cap) are carried by the collective primitives in
\texttt{quatrex.core.mpi\_linalg} (\texttt{global\_dot}, \texttt{global\_norm},
\texttt{global\_gram}, \texttt{trust\_cap}), so the serial linear algebra written
below is evaluated across the distributed energy grid.

\paragraph{Reduced-rank extrapolation (RRE).}
RRE~\cite{sidi_vector_2017} forms the affine combination
$\Sigma=\sum_i\gamma_i\,\Sigma^{(i)}$, $\sum_i\gamma_i=1$, of a short window of
iterates that minimises the residual norm
$\bigl\|\sum_i\gamma_i\,\vc{f}_i\bigr\|$ (a minimal-residual vector extrapolation;
distinct from minimal-polynomial extrapolation, MPE).
Restarted every few iterates, it recovers a mildly unstable fixed point from the
algebra of the residual history rather than by contraction, at the cost of one
small real Gram solve per cycle and no extra map evaluations.

\paragraph{Recursive projection method (RPM).}
When the unstable subspace is low-dimensional, the RPM of
\citet{shroff_stabilization_1993} applies Newton on that subspace and Picard on
its contractive complement,
% REVIEW: corrected to the canonical Shroff-Keller RPM update. The correction
% term is +Z(I_k-H)^{-1} H Z^H f, whose modal gain on an unstable mode is
% lambda/(1-lambda) (matching phi(lambda) in the text below); the previous
% -Z(I_k-H)^{-1} Z^H f dropped the H factor and carried the wrong sign.
\begin{equation}\label{eq:rpm}
  \Sigma^{(n+1)}
  =\mathcal{F}[\Sigma^{(n)}]
   +Z\,(\mathbb{I}_k-H)^{-1} H\,Z^{H}\vc{f}^{(n)},
  \qquad \vc{f}^{(n)}=\mathcal{F}[\Sigma^{(n)}]-\Sigma^{(n)},
\end{equation}
where the columns of $Z$ span the $k$-dimensional unstable invariant subspace and
$H=Z^{H}G'Z$ is the restricted Jacobian. Both are obtained without extra map
evaluations -- $Z$ from the dominant singular vectors of a sliding buffer of
increments (a dynamic-mode decomposition) and $H$ from the projected secants
$Z^{H}\Delta\mathcal{F}=H\,Z^{H}\Delta\Sigma$. The correction multiplies
$\vc{f}$ and so vanishes identically at the fixed point; the safeguarded modal
step $\phi(\lambda)=\lambda/(1-\lambda)$ is applied only on directions with
$|\lambda(H)|>1$, the rest being left to Picard.

\paragraph{Jacobian-free Newton--Krylov (JFNK).}
When the unstable subspace is larger, or its eigenvalues very large, we instead
solve the Newton system for the residual $R(\Sigma)=\mathcal{F}[\Sigma]-\Sigma$,
\begin{equation}\label{eq:jfnk}
  (G'-\mathbb{I})\,\delta=-R,
\end{equation}
by GMRES~\cite{saad_gmres_1986}, which requires only Jacobian--vector products.
These are formed matrix-free by a finite difference of the map,
\begin{equation}\label{eq:jfnk_jv}
  (G'-\mathbb{I})\,\vc{v}
  =\frac{\mathcal{F}[\Sigma+\epsilon\vc{v}]-\mathcal{F}[\Sigma]}{\epsilon}-\vc{v},
\end{equation}
at one self-energy evaluation per Krylov vector. Because the self-energy map
conjugates -- $\Sigma^R$ through the Kramers--Kronig relation
\cref{eq:sigma_R_full}, $G^A=(G^R)^{H}$, and the bubble leg structure -- $G'$ is
real-linear but not complex-analytic, so the Krylov solve is run in the real
embedding $[\,\mathrm{Re}\,\Sigma,\ \mathrm{Im}\,\Sigma\,]$. The inner solve is
truncated inexactly with an Eisenstat--Walker forcing
term~\cite{eisenstat_choosing_1996} and globalised by a trust-region cap on
$\|\delta\|$, following~\citet{brown_hybrid_1990,kelley_iterative_1995}. A
Levenberg--Marquardt / pseudo-transient shift
$G'-\mathbb{I}\to G'-\mathbb{I}+\mu\mathbb{I}$ is inappropriate here: the Newton
Jacobian of \cref{eq:jfnk} is indefinite (contractive modes contribute
eigenvalues in $[-2,0)$, unstable modes positive eigenvalues), so a positive
shift drives a negative eigenvalue through zero; the trust region is the only
safe globalisation.

\subsubsection{The $\eta\to0$ limit and grid-consistent regularisation}
\label{sub:eta_zero}

The broadening $\eta$ in \cref{eq:system} is a numerical device; the conserving
transport limit is $\eta\to0$. Its difficulty is a resolvent-sensitivity effect.
Differentiating the Dyson equation gives
$\mathrm{d}G^R=G^R(\mathrm{d}\Sigma^R)G^R$, so near a mode of frequency $\Omega$
and linewidth $\Gamma$ the map gain scales as
\begin{equation}\label{eq:resolvent_gain}
  \Bigl|\frac{\partial G^R}{\partial\Sigma^R}\Bigr|\sim|G^R|^2\sim
  \frac{1}{(2\Omega\Gamma)^2},
  \qquad
  |\lambda(G')|\sim\frac{\Phi^2}{\max(\Gamma,\,c\,\Delta\omega)^2},
\end{equation}
the floor $c\,\Delta\omega$ arising because a Lorentzian narrower than the grid
samples as an undamped pole. The iteration is therefore marginal or unstable
wherever a heat-carrying mode is sharper than the grid,
$\Gamma\lesssim\Delta\omega$ -- the regime the root-finders above are built for.
We pair them with two regularisations that vanish as $\Delta\omega\to0$ and
preserve the $\Phi$-derivable conservation laws.

\paragraph{Smooth spectral window.}
Above the harmonic band top and inside any spectral gap the device Green's
functions have no support, yet the convolution \cref{eq:sigma_phph} leaks a small
self-energy there which, at $\eta\to0$, is divided by an undamped resolvent.
Masking those bins with a hard window is itself harmful: a step discontinuity in
$\Sigma$ has a Kramers--Kronig partner $\sim\ln|\omega-\omega_{\rm edge}|$ that
rings (Gibbs), manufacturing a marginal mode at the mask edge. We instead apply a
\emph{smooth} multiplicative window
$w(\omega)=w_{\rm supp}(\omega)\,w_{\rm occ}(\omega)\in[0,1]$: a raised-cosine
ramp of width a few $\Delta\omega$ on the distance to the harmonic support
(computed once from the band structure $D(k)=D_0+D_1e^{ik}+D_1^{H}e^{-ik}$,
independent of $G$), times a logistic in the Bose occupation that freezes
thermally inactive spectator modes. The window multiplies the input
Green's-function legs, the output $\Sigma^{\gtrless}$, and the Kramers--Kronig
input; because the same per-frequency factor multiplies both sides of the
energy-balance identity \cref{eq:bubble_balance}, the equality
$P_{\rm in}=P_{\rm out}$ is preserved and conservation is left intact.

\paragraph{Infrared occupancy taper.}
The Bose factor $n(\omega)\simeq k_BT/\hbar\omega$ diverges as $\omega\to0$;
sampled at the first grid bin it injects a spurious $\sim 1/\Delta\omega$
linewidth, whereas the acoustic sum rule forces $\Gamma\sim\omega^2\to0$. We
restore the correct onset with the taper
\begin{equation}\label{eq:ir_taper}
  t(\omega)=\frac{\omega^2}{\omega^2+\omega_{\rm reg}^2},
  \qquad\omega_{\rm reg}=C\,\Delta\omega,
\end{equation}
applied identically to both contact occupations so that every $G^<$ leg inherits
it and the balance is preserved; as $\Delta\omega\to0$, $\omega_{\rm reg}\to0$ and
$t\to1$, and the physical conductance is the $\omega_{\rm reg}\to0$ extrapolant.
This taper is a \emph{pragmatic} regulator: it is a cutoff, removed by the
$\omega_{\rm reg}\to0$ extrapolation, and because it multiplies the occupations it
slightly suppresses the genuine low-frequency transport. That suppression is small
--- the heat current is $\hbar\omega$-weighted, so the infrared region carries
little conductance --- and, as \cref{sub:ir_subtraction} shows, the bare untapered
bubble is itself exactly conserving; the taper is therefore a stabiliser of the
marginal $\eta=0$ iteration, not a correction to the physics.

\paragraph{Frequency-proportional regulator and the sub-grid soft-mode floor.}
The constant broadening $(\omega^2+i\eta)\mathbb{I}$ of \cref{eq:system} is a
convenient idealisation; the production solver instead uses a
\emph{frequency-proportional} regulator
\begin{equation}\label{eq:regulator_impl}
  [G^R]^{-1} = \bigl(\omega^2 + 2i\eta|\omega|\bigr)\mathbb{I} - D - \Sigma^R,
\end{equation}
i.e.\ $\omega^2$ shifted by $2i\eta|\omega|$ rather than by $i\eta$ or
$(\omega+i\eta)^2$. The $(\omega+i\eta)^2$ form carries a spurious $-\eta^2$
band-edge shift that makes modes with $\omega\lesssim\eta$ artificially evanescent
and suppresses the acoustic plateau; the linear-in-$|\omega|$ damping avoids this
and vanishes as $\omega\to0$, so it does not pollute the infrared. The price is
that the acoustic \emph{soft} modes ($D\to0$ as $\omega\to0$) are then left
unregularised at $\eta\to0$ ($G^R\sim1/\omega^2$) and the marginal SCBA diverges.
We restore conditioning with a DC-concentrated soft-mode broadening floor added to
the \emph{device} system matrix only (never to the ideal-lead OBC),
\begin{equation}\label{eq:eta_floor}
  \Delta z^2(\omega) = i\,\Gamma_{\rm floor}\,\frac{\omega_c^2}{\omega^2+\omega_c^2},
  \qquad \Gamma_{\rm floor}=(c_{\rm floor}\,\Delta\omega)^2,\quad
  \omega_c = 2\Delta\omega,
\end{equation}
a Lorentzian low-pass that damps only the lowest, grid-unresolved (and essentially
heat-free) bins while leaving the resolved low-$\omega$ physics and the plateau
untouched. Since $\Gamma_{\rm floor}\propto(\Delta\omega)^2\to0$ under grid
refinement, it is grid-consistent and removed in the same limit as the taper
\cref{eq:ir_taper}. The floor is \emph{annealed} during the SCBA: $c_{\rm floor}$
ramps from a larger start to its target over the first iterations
(the config knobs \texttt{eta\_ir\_floor\_cells}~$\to$~%
\texttt{eta\_ir\_floor\_final\_cells} over
\texttt{eta\_ir\_floor\_ramp\_iterations}), so the soft modes are damped while
$\Sigma$ develops and the broadening is then relaxed toward zero at the converged
fixed point. This floor-plus-anneal --- not the spectral window or the occupancy
taper alone --- is what makes the $\eta=0$ fixed point of soft systems such as the
\texttt{d5a} nanowire reachable (\cref{sec:res_eta0}).

\paragraph{Linewidth feasibility criterion.}
Whether $\eta=0$ is reachable on a given grid can be decided a~priori, without any
SCBA, by evaluating the golden-rule linewidth $\Gamma_{\rm anh}(\omega)$
(\cref{eq:fgr}) directly from $\Phi^{(3)}$ and the harmonic eigenvectors and
comparing it to $\Delta\omega$: by \cref{eq:resolvent_gain} the fixed point is
well posed only where every heat-carrying mode satisfies
$\Gamma_{\rm anh}\gtrsim\Delta\omega$.

\paragraph{Energy-balance diagnostic.}
Beyond the lead imbalance $\delta I/\bar I$, full $\Phi$-derivable
self-consistency implies the per-iteration identity
\begin{equation}\label{eq:bubble_balance}
  P_{\rm in}=\sum_\omega\hbar\omega\,
   \mathrm{Tr}\bigl[\Sigma^<(\omega)G^>(\omega)\bigr]
  =\sum_\omega\hbar\omega\,
   \mathrm{Tr}\bigl[\Sigma^>(\omega)G^<(\omega)\bigr]=P_{\rm out},
\end{equation}
the net power exchanged by the scattering self-energy. Its residual isolates
vertex and convolution errors from contact errors and is the most stringent of
our convergence checks.

\subsubsection{Infrared limit of the heat current}
\label{sub:ir_subtraction}

The taper \cref{eq:ir_taper} regularises the iteration, not the physics; the
infrared behaviour itself is benign, and no exact bubble-level subtraction is
required. The Bose factor $n(\omega)=1/(e^{\beta\hbar\omega}-1)$ enters three
places --- the contact injection
$\Sigma_\ell^{\lessgtr}=i\Gamma_\ell[n_\ell+(1{\pm}1)/2]$, the bubble legs
$G^{\lessgtr}(\omega')\sim n(\omega')A(\omega')$ of \cref{eq:sigma_guo}, and the
Meir--Wingreen integral \cref{eq:meir_wingreen} --- through its Laurent (Bernoulli)
expansion
\begin{equation}\label{eq:bose_laurent}
  n(\omega,T)=\frac{k_BT}{\hbar\omega}-\frac12+\frac{\hbar\omega}{12k_BT}
  -\frac{(\hbar\omega)^3}{720(k_BT)^3}+\cdots,
\end{equation}
whose singular part is the c-number pole $k_BT/\hbar\omega$.

\paragraph{The pole is physical only in the observable.}
In the heat current the $1/\omega$ is cancelled by the explicit $\hbar\omega$
weight. The constant $-\tfrac12$ is temperature-independent and drops from any lead
difference, $\hbar\omega\,[n_L-n_R]=k_B(T_L-T_R)+\mathcal{O}(\omega^2)$, so the
spectral density has the finite infrared limit
\begin{equation}\label{eq:ir_plateau}
  j(\omega)=\frac{\hbar\omega}{2\pi}\,\mathcal{T}(\omega)\,[n_L(\omega)-n_R(\omega)]
  \xrightarrow[\omega\to0]{}\frac{k_B\,\Delta T}{2\pi}\,\mathcal{T}(0),
  \qquad \mathcal{T}(0)=N_{\rm ac},
\end{equation}
with $N_{\rm ac}$ the number of acoustic channels, each contributing the quantised
thermal conductance $g_0=\pi^2k_B^2T/3h$~\cite{rego_quantized_1998,schwab_measurement_2000,jeong_analysis_2012}
(\cref{app:infrared} for the derivation). This plateau is the physical infrared
target; a regulator that multiplies the occupations by $\omega^2$, as the taper
\cref{eq:ir_taper} does, removes a finite, quantised piece of it
(\cref{fig:res_eta0_ir}).

\paragraph{There is no pole to subtract in the bubble.}
It is tempting to remove the $1/\omega$ from the scattering self-energy by the
singularity-subtraction device of boundary-integral
numerics~\cite{helsing_higher_2013,hanninen_singularity_2006,wang_logarithmic_2023}.
Applied to the bubble legs this is both unnecessary and harmful.
% REVIEW: removed "|V_3|^2 propto omega_1 omega_2 omega_3, so the soft-leg factor
% offsets the absorption pole". In the harmonic mode basis each leg carries
% sqrt(hbar/2 omega) ~ omega^{-1/2}, so |V_3|^2 propto 1/(omega_1 omega_2 omega_3),
% which COMPOUNDS the n(omega')~1/omega' pole rather than offsetting it. Finiteness
% comes from the linearly vanishing acoustic spectral weight (the same omega-factors,
% carried by A) together with the acoustic phase space (~omega'^2 in 3D).
The acoustic spectral weight vanishes
linearly, $A(\omega')\to A'(0)\,\omega'$, so the equilibrium lesser function
$G^{<}_{\rm eq}=-i\,nA$ stays finite despite $n(\omega')\sim1/\omega'$. Out of equilibrium the bosonic fold
$G^{<}(-\omega)=G^{>}(\omega)^T$ together with the bounded spectral function forces
the $1/\omega$ parts of $G^{<}$ and $G^{>}$ to be equal and opposite, hence both
zero; and the contact broadening $\Gamma_\ell(\omega)$, being the bosonic (odd)
anti-Hermitian part of $\Sigma^R_\ell$, vanishes as $\omega\to0$, so
$\Gamma_\ell n_\ell$ is finite. The device bubble therefore carries no net infrared
pole. Subtracting one regardless feeds the bubble a modified $G$, so that
$\Sigma\neq\mathcal{F}[G]$ and the $\Phi$-derivable balance \cref{eq:bubble_balance}
is broken. The accurate and conserving choice near $\omega=0$ is the \emph{bare}
bubble; the taper is retained only as a numerical stabiliser of the marginal
$\eta=0$ iteration (\cref{sub:eta_zero}), its small effect on the conductance
removed by the $\omega_{\rm reg}\to0$ extrapolation rather than by any analytic
subtraction. In a single nanotube the long-wavelength three-phonon linewidths scale
per polarisation ($\Gamma_{\rm LA}\sim|k|^{-1/2}$, $\Gamma_{\rm flexural}\sim|k|^{0}$,
$\Gamma_{\rm twist}\sim|k|^{1/2}$)~\cite{bruns_nanotube_2021,ma_examining_2014,maznev_demystifying_2014},
so the converged anharmonic plateau need not equal the ballistic
$\mathcal{T}(0)=N_{\rm ac}$.

\subsubsection{Convergence diagnostics}

We monitor three quantities across iterations: (i) the relative
Frobenius norm
$\|\Sigma_s^{(n+1)}-\Sigma_s^{(n)}\|/\|\Sigma_s^{(n)}\|$,
(ii) the conductance $G(T)$ from \cref{eq:G_thermal} or, in the
inelastic case, the Meir--Wingreen current, and (iii) the slab-resolved
current imbalance $\delta I/\bar I$. Convergence is declared when (i)
and (iii) fall below prescribed tolerances simultaneously. Under a
temperature bias at finite $\eta$, however, (iii) does not reach zero but
floors at the ghost-reservoir term of \cref{eq:sumrule}
(\cref{app:conservation}); the primary gates are then (i) and the
pairing-exact bubble balance $P_{\rm in}=P_{\rm out}$, with (iii) reserved as
an $\eta$-budget diagnostic.

\subsubsection{Lowest-order approximation.}

% REVIEW: re-attributed. The one-shot current-conserving lowest-order
% approximation is due to Mera et al.; the lee_* papers apply/extend it to
% anharmonic phonon transport in 3D nanostructures.
An attractive alternative to direct SCBA is the
lowest-order approximation (LOA), introduced for inelastic quantum transport
by~\citet{mera_inelastic_2012} and applied to anharmonic phonon transport
by~\citet{lee_anharmonic_2018,lee_quantum_2019}, implemented through Padé
resummation. LOA collects only the conserving ($\Phi$-derivable) diagrams
at each order in $\Phi^{(3)}$, ensuring exact current conservation
at every iteration, then accelerates the resulting series by
Padé approximants. Compared to SCBA, LOA achieves comparable accuracy
at $\sim 2\times$ lower cost for thermal transport in 3D
nanostructures~\cite{lee_anharmonic_2018}.

% ----------------------------------------------------------------------
\subsection{Decay properties and Truncation errors}
\label{sub:decay}
% ----------------------------------------------------------------------

Truncations at every step (real-space FC cutoff, finite device length, finite $\omega$ grid,
finite Sancho-Rubio iteration count) introduce errors with
quantifiable scaling. We summarize, from first principles, the decay
rates of each of the involved quantities.

\subsubsection{Real-space decay of harmonic force constants}
\label{subsub:fc2_decay}

The asymptotic behaviour of $\Phi^{(2)}(R)$ at large $|R|$ is controlled
by the analytic structure of the dielectric response of the underlying
electronic system. \Citet{cances_decay_2025} prove the following
within the reduced Hartree--Fock (or random-phase
approximation) model:

\paragraph{Insulators at zero electronic temperature.}
The dynamical matrix $D(\qp)$ is singular at $\qp = 0$, with the
non-analytic dipole--dipole structure given in~\cite{cochran_dielectric_1962}:
\begin{equation}\label{eq:fc2_dipole}
  D^{\mathrm{LR}}_{\kappa\alpha,\kappa'\beta}(\qp \to 0)
  \;\sim\;
  \frac{4\pi}{\Omega_{\mathrm{cell}}}\,
  \frac{(\qp\cdot Z^*_{\kappa})_\alpha\,(\qp\cdot Z^*_{\kappa'})_\beta}
       {\qp\cdot \epsilon_\infty\cdot \qp},
\end{equation}
where $Z^*_\kappa$ is the Born effective charge tensor of atom $\kappa$
and $\epsilon_\infty$ the optical dielectric tensor. Inverse Fourier
transform of this $\qp\to 0$ singularity gives an algebraic
real-space tail to $\Phi$.

\paragraph{Materials at finite electronic temperature (metals).}
\Citet{cances_decay_2025} prove that the
dynamical matrix is analytic in $\qp$, and
$\Phi^{(2)}(R)$ decays exponentially:
\begin{equation}\label{eq:fc2_metal}
  |\Phi^{(2)}(R)|\;\le\;C\,e^{-|R|/\xi_e},\qquad
  \xi_e \;\sim\;\frac{\hbar v_F}{k_B T_{\mathrm{el}}}.
\end{equation}

\paragraph{Concrete truncation error.}
Truncating at distance $R_c$ introduces an error in the dynamical matrix
of order
\begin{equation}\label{eq:fc2_error_metal}
  \|D(\qp) - D^{R_c}(\qp)\| \;\lesssim\;C \frac{e^{-R_c/\xi_e}}{1 - e^{-1/\xi_e}}
  \quad\text{(metals, exponential)},
\end{equation}
or, in polar insulators, $\sim 1/R_c^3$ for the bare truncation, but
exponential, $\sim e^{-R_c/\xi_g}$, once the
$1/r^3$ Cochran--Cowley tail has been subtracted analytically and added
back exactly via Ewald summation~\cite{gonze_dynamical_1997}.
% REVIEW: the insulator short-range decay length xi_g is set by the electronic
% band gap, NOT by the metal thermal length xi_e = hbar v_F/(k_B T_el) of
% eq:fc2_metal (an insulator at T=0 has no Fermi velocity). Also the geometric
% tail in eq:fc2_error_metal is the 1D form; in d dimensions a shell degeneracy
% multiplies the summand by ~R_c^{d-1}.

\subsubsection{Transverse $\qp$-mesh convergence}
\label{subsub:qp_decay}

The convergence of observables in $N_q$ depends on the differentiability
of the integrand. For the Caroli formula \cref{eq:caroli}, the
$\qp$-integrand is generically smooth and trapezoidal Monkhorst--Pack
sampling converges as $\mathcal{O}(N_q^{-2})$ in 2D. At Brillouin-zone boundaries or near van
Hove singularities, convergence degrades to $\mathcal{O}(N_q^{-1})$;
a Gaussian smearing or tetrahedron method may then be preferable. The
self-energy convolution \cref{eq:sigma_phph} couples $\qp$ and $\qp{}'$,
so the effective integrand has $N_q^{-1/2}$ statistical-noise structure
% REVIEW: importance sampling lowers the variance prefactor but not the
% N_q^{-1/2} Monte-Carlo exponent; approaching O(N_q^{-1}) needs quasi-Monte-Carlo
% / low-discrepancy sequences.
under random sampling, and reducing this toward $\mathcal{O}(N_q^{-1})$ requires
variance reduction beyond plain importance sampling, e.g.\ quasi-Monte-Carlo /
low-discrepancy sampling~\cite{lee_anharmonic_2018}.

\subsubsection{Convergence rate of Sancho-Rubio decimation revisited}

As discussed in \cref{sub:contact_se}, the decimation iteration
\cref{eq:sr_eps_s,eq:sr_eps,eq:sr_alpha,eq:sr_beta} effectively doubles
the number of decimated layers at each step, with iterates obeying
$\|\alpha_n\|, \|\beta_n\| \sim \rho^{2^n}$ inside gaps and
$e^{-c\eta\cdot 2^n}$ on bands. The practical implication is that
\emph{doubling the tolerance roughly costs one extra iteration}; total
cost is $\mathcal{O}(N_\omega\,N_q\,N_{\mathrm{DOF}}^3\,N_{\mathrm{iter}})$
with $N_{\mathrm{iter}}$ essentially constant (and small,
$\sim 10\text{--}30$).

% ----------------------------------------------------------------------
